\section{The Demonstration Playground}\label{sec:implementation-playground}

While the core contribution of this thesis is the compiler, a web-based \textit{Playground} was developed to demonstrate the language's capabilities and provide an intuitive environment for composition.
The playground integrates the C++ compiler into a modern full-stack web application.

\subsection{System Architecture}

The playground follows a distributed architecture:
\begin{itemize}
    \item \textbf{Frontend (Next.js):} A React-based interface that provides the editor and visualization tools.
    \item \textbf{Backend (Express):} A Node.js API that acts as a bridge. It receives DSL source code from the client, invokes the C++ \texttt{dslrc} binary via a sub-process, and streams the resulting MIDI file back to the browser.
    \item \textbf{Integration:} The system is containerized using Docker, ensuring that the C++ environment (including Bison/Flex runtimes) is identical across development and production.
\end{itemize}

\subsection{IDE Integration and Diagnostics}

The playground features an IDE-like experience powered by the Monaco Editor.
A custom "Monarch" tokenizer was implemented for the DSL to provide syntax highlighting for musical keywords, pitch literals, and control flow.

The most significant integration is the \textit{Diagnostics Feedback Loop}.
When a compilation fails, the Express server parses the structured error output from the compiler and returns it to the client.
The React application then converts these diagnostics into Monaco "Markers," which underline the erroneous code and display the compiler's semantic or syntax error messages in a tooltip.
This real-time feedback is enabled by the "insane" compilation speeds discussed in the evaluation chapter.

\subsection{The Piano Roll and Synthesis}

To visualize the composition, a custom \textit{Piano Roll} component was built using the HTML5 Canvas API\@.
Upon a successful compilation, the browser parses the returned MIDI bytes and renders them as colored rectangles on a scrolling timeline.

The synthesis engine utilizes the \textit{Tone.js} library.
Because MIDI files contain only instructions and no sound, the playground loads a set of high-quality \textit{Soundfonts} representing General MIDI instruments.
The playback logic schedules notes in the browser's audio context, allowing the user to hear their "coded" composition instantly.
The synchronization between the scrolling piano roll and the audio playback is maintained through a high-precision request-animation-frame loop, ensuring a professional and "alive" feel to the prototype.
